\documentclass[aspectratio=169]{beamer}
\usetheme{Madrid}
\setbeamertemplate{navigation symbols}{}
\setbeamertemplate{footline}[page number]

\newcommand{\Title}{CPA using the resolvent formalism}
\title{\Title}
\author{Frank Takara Ebel}
\institute{
    \begin{columns}
        \column{0.3\textwidth}
        Student ID\@: 01429282 \\
        Programme Code: UE 796 700 461
    \end{columns}
}
\date{} % no date

\input{preamble/packages.tex}
\input{preamble/commands.tex}

\begin{document}

\frame{\titlepage}

\section*{main}

\begin{frame}{Motivation}
    In exercise used finite broadening $\delta$:
    \begin{align}
        G_{0}(k, \omega+\mi\delta) & = \frac{1}{\omega + \mi\delta - \epsilon_k} \\
    \end{align}
    Can we do $\lim_{\delta\to0^+}$? \\
    Yes, in the resolvent formalism:
    \begin{equation}
        \res(\omega, H) =  (\omega\mathbb{1} - H)^{-1}
    \end{equation}
\end{frame}

\end{document}
